\section{The SPIKE Framework}
\label{sec:method}

% 中文：我们将长时程开放世界控制建模为一个包含多模态观测、分层记忆、显式恢复和 token--latency 预算的序列决策问题。
% 在每一步，智能体根据当前视觉场景、文本上下文、推断的 UI 状态、近期动作和可选执行日志选择可执行动作。
We model long-horizon open-world control as sequential decision making with
multimodal observations, layered memory, recovery, and token--latency budgets:
the Event Trigger decides when to deliberate, the Reactive Controller acts
between deliberations, and Hierarchical Memory retrieves evidence at routine and
escalation boundaries.


Open-world control has a cost-value mismatch: routine steps need fast execution,
whereas rare scene changes, stalls, and failures require expensive reasoning.
Always-strategic agents spend large-model calls on routine motion; purely
reactive agents lack reliable triggers for stale plans. \method{} treats
strategic reasoning as event-triggered, amortizes proposals over stable segments,
and retrieves different evidence for local and escalated steps.

\begin{figure}[t!]
    \centering
    \includegraphics[width=0.98\linewidth]{Architecture.pdf}
    \caption{\textbf{Closed-loop event-triggered deliberation in \method{}.} The Strategic Controller performs information collection,
    optional Self Reflection, task reasoning, and short-horizon action proposal;
    the Reactive Controller uses real-time signals and retrieved memory to
    follow, override, or trigger replanning. The Hierarchical Memory module
    combines Global Memory, the State-Action Memory Bank (\samb{}), and the
    State-Action Knowledge Graph (SA-KG) through vector search and graph
    traversal, then propagates updates after action execution.}
    \label{fig:overview}
    \vspace{-0.5em}
\end{figure}

\subsection{Event-Triggered Dual Controller Framework}
\label{subsec:overview}

% 中文：受双过程认知理论启发，\method{} 采用非对称双控制器设计。
% \strategicsystem{} 负责低频的全局规划、失败分析和任务级重定向；\reactivesystem{} 在严格 token 预算下负责低延迟局部执行。
% 两个控制器由 Event Trigger 协调；Hierarchical Memory 为不同控制器提供各自所需的上下文。
For an episode with horizon $H$, let
$o_t=(v_t,u_t,f_t)$ denote the observation at step $t$, where $v_t$ is the
visual screenshot, $u_t$ is textual or user interface (UI) context, and $f_t$ is runner feedback
such as progress, position, inventory, or dialogue changes. The agent chooses an
executable game action $a_t\in\mathcal{A}$ from history
$h_t=(o_1,a_1,\ldots,o_t)$. An event-trigger decision $g_t\in\{0,1\}$ selects the
Strategic Controller ($g_t=1$) or Reactive Controller ($g_t=0$). For trajectory
$\tau=(o_1,a_1,\ldots,o_H)$, success is $S(\tau)\in\{0,1\}$ and the deliberation
budget is $B$. Let $C_t(g_t,h_t)$ denote the history-dependent model-call, token,
or latency cost. The target is
\begin{equation}
\label{eq:budgeted_control_objective}
\max_{\pi_R,\pi_S,g}\mathbb{E}[S(\tau)]
\quad
\mathrm{s.t.}\quad
\mathbb{E}\!\left[\sum_{t=1}^{H} C_t(g_t,h_t)\right]\le B .
\end{equation}
Inspired by dual-process theories of cognition \citep{kahneman2011thinking},
\method{} separates slow, high-cost planning from fast local execution. The
Strategic Controller handles planning, failure analysis, and task redirection;
the Reactive Controller executes local actions under a strict token budget; the
Event Trigger and Hierarchical Memory allocate deliberation and evidence.
This differs from always-deliberate agents, which call the expensive planner at
nearly every step, from fixed fast-slow pipelines, which switch controllers by a
static schedule, and from memory-augmented agents with a single retrieval pool.
In \method{}, controller allocation, local override, and memory retrieval are
all conditioned on detected event boundaries, so the framework changes both
\emph{when} reasoning is invoked and \emph{what} evidence each controller sees.

% 中文：如图~\mbox{\ref{fig:overview}} 所示，\method{} 包含 Env and Information Capture、Strategic Controller、Hierarchical Memory、Reactive Controller 和 Event Trigger 五个模块。
% 系统遵循统一闭环：信息捕获、路由、记忆检索、动作提议、动作执行与记忆更新。
% Env and Information Capture 提供实时视觉和文本信息并刷新 Hierarchical Memory；Event Trigger 检查场景变化、卡死状态和进展信号。
% 常规状态由 \reactivesystem{} 根据 State-Action Memory Bank 中的 Memory Information 和当前子任务直接处理。
% 当检测到突发视觉变化、重复失败或进展不足时，Event Trigger 将控制权升级给 \strategicsystem{}，生成未来 $T$ 步的短程计划。
% 随后 \reactivesystem{} 可以遵循该计划、在局部证据表明计划失效时覆盖它，或请求重规划；动作执行后，Memory Update 将结果写回局部和长期记忆。

As illustrated in Figure~\mbox{\ref{fig:overview}}, \method{} consists of Env
and Information Capture, Strategic Controller, Hierarchical Memory, Reactive
Controller, and Event Trigger. Together they form a closed loop:
capture produces $o_t$ and updates $h_t$; the Event Trigger scores whether the
current history marks an event boundary; the selected controller retrieves
controller-specific memory and emits $a_t$; execution feedback updates
\samb{}, SA-KG, and $h_{t+1}$; the next step repeats this trigger--act--update
cycle.
\vspace{-0.35em}

\paragraph{Controller roles.}
The Strategic Controller is activated when periodic refresh, stuck-state, or
recovery signals require deliberate reasoning. It integrates the current
observation, task summary, recent failure traces, and structured long-term
memory, then follows the workflow
\begin{center}
\begingroup
\vspace{-0.35em}

\begin{tikzpicture}
\node[
    draw=boxteal,
    fill=oursrowblue,
    rounded corners=5pt,
    line width=0.8pt,
    inner xsep=6pt,
    inner ysep=4pt,
    text width=0.88\linewidth,
    align=center
] {Information Collection -> Self Reflection -> Task Reasoning -> Action Proposal.};
\end{tikzpicture}
\vspace{-0.35em}

\endgroup
\end{center}
Figure~\mbox{\ref{fig:workflow}} expands the Strategic Controller workflow. The
optional Self Reflection stage handles unclear evidence, contradiction, or
repeated failure. The output is a short-horizon proposal for the next $T$ steps,
long enough to guide local execution and compact enough to revise after a scene
change.

The Reactive Controller handles routine steps under a strict token budget. It
uses the latest strategic proposal, current observation, and Memory Information
from \samb{} to emit local actions. When new evidence invalidates the proposal,
it may only apply bounded local corrections such as avoidance, repositioning, or
tool reselection; it cannot change the active subgoal, and failed corrections
return control to the Event Trigger.

\begin{figure}[t!]
    \centering
    \includegraphics[width=0.94\linewidth]{workflow.pdf}
    \caption{\textbf{Strategic Controller workflow in one escalated planning cycle.}
    The example shows how \method{} collects scene evidence, checks progress,
    reasons over the task, and emits a short plan for Reactive Controller execution.}
    \label{fig:workflow}
    \vspace{-0.5em}
\end{figure}

\subsection{Event Trigger as an Online Surrogate for Budgeted Deliberation}
\label{subsec:controller_scheduler}

The Event Trigger operationalizes the budgeted objective as a practical online
event-boundary decision. Ideally, each step would choose whether the expected
utility of strategic deliberation justifies its extra cost:
\vspace{-0.35em}
\begin{equation}
\label{eq:event_trigger_optimization}
g_t^\star=\arg\max_{g\in\{0,1\}}
\widehat{U}(g\mid h_t)-\lambda C_t(g,h_t),
\end{equation}
\vspace{-0.1em}
where $\widehat{U}(g\mid h_t)$ estimates progress and recovery value from
history and $\lambda$ controls the cost penalty. Because the true value of
deliberation is not observable online from sparse task success, \method{} uses a
reproducible surrogate event detector. Eq.~\mbox{\ref{eq:scheduler_routing}} uses
fixed thresholds over visual discontinuity, progress stagnation, behavioral
repetition, and execution failure; these signals identify stale strategic
proposals or unreliable local execution rather than estimating task success.
Appendix Sec.~\mbox{\ref{sec:appendix_event_algorithm}} gives the trigger--act--update loop.

The trigger variable $g_t \in \{0,1\}$ indicates whether step $t$ escalates:
\vspace{-0.35em}
\begin{equation}
\label{eq:scheduler_routing}
g_t=\mathbb{I}\!\left[(c_t\ge T)\vee(d_t>\tau_v)\vee(z_t\ge\tau_z)\vee(r_t\ge\tau_r\wedge z_t\ge\tau_{rz})\vee(\ell_t\ge\tau_\ell)\right].
\end{equation}
where $c_t$ counts Reactive steps since the last Strategic call and resets after
it, $T$ is the refresh interval, and
$\tau_v,\tau_z,\tau_r,\tau_{rz},\tau_\ell$ are fixed triggering thresholds.
The detector maps signals to observable failure modes: $d_t=1-\cos(\phi(v_t),\phi(v_{t-1}))$
catches stale plans after visual discontinuities, $z_t$ catches stalled routines,
$r_t$ catches behavior loops within the last $W$ actions, and
$\ell_t\in\{0,1,2,3\}$ catches execution failure. We set
$\Delta_t=q_t-q_{t-1}$, where $q_t$ is a deterministic progress counter derived
from the shared \stardojo{} runner feedback, not a privileged game-state API;
$z_t=0$ when $\Delta_t>0$ and $z_t=z_{t-1}+1$ otherwise. All baselines receive
the same raw observation, action schema, execution feedback, and history fields;
\method{} only adds Event Trigger bookkeeping.

The fixed thresholds instantiate the event detector rather than a per-task
policy; Table~\mbox{\ref{tab:scheduler_sensitivity}} reports their local
sensitivity, and Appendix Sec.~\mbox{\ref{sec:appendix_scheduler_details}}
details their computation. These thresholds are tied to event types rather than
individual tasks: within the same game family, where camera motion, action
granularity, and progress feedback have similar scales, the same settings can be
reused without task-by-task replacement. When moving to a different game type,
the Event Trigger should be recalibrated according to the new event statistics,
such as visual dynamics, interaction pace, stall duration, and failure-message
density.
Failure levels scope recovery: F1 keeps bounded retry local, while F2/F3
escalate invalid actions, stalls, or repeated compound failures when local
context becomes unreliable.

\FloatBarrier
\subsection{Hierarchical Memory}
\label{subsec:memory}

% 中文：3.3 只保留记忆层的核心主线和两个检索打分公式，字段、融合和写回细节放到附录。
Long-horizon tasks risk context overflow, hallucination, and forgetting.
Figure~\mbox{\ref{fig:overview}}(c) organizes Hierarchical Memory around
controller-specific roles: Global Memory preserves task continuity, the
State-Action Memory Bank (\samb{}) is a low-latency local action prior for
Reactive execution, and the State-Action Knowledge Graph (SA-KG) is a strategic
evidence graph for checking and replanning rather than a flat retrieval store.

Global Memory keeps a fixed-size window of recent observations and actions,
together with the active subgoal, progress indicators, and failure summaries.
When the Event Trigger keeps execution local, the Reactive Controller
retrieves nearby state-action traces from \samb{} as local hints, with the risk
of stale habits reduced by recency, reward, and success terms. Given the current
query $x_t$, \samb{} ranks item $i$ by
\begin{equation}
\label{eq:samb_score}
S_{\mathrm{MB}}(i \mid x_t)=\rho_i(\alpha_C C_i+\alpha_L L_i+\alpha_R R_i+\alpha_P P_i).
\end{equation}
All terms are normalized before weighting: $L_i$ is the Jaccard overlap between
lowercased query and memory token sets, $C_i$ is token-frequency cosine
similarity, $R_i=(\mathrm{clip}(\mathrm{rewardEMA}_i,-1,1)+1)/2$ uses an
exponential moving average (EMA) of observed rewards, $P_i$ is
successes over attempts, $\rho_i=\exp(-\mathrm{ageHours}_i/24)$, and
$\alpha_C,\alpha_L,\alpha_R,\alpha_P$ are fixed retrieval weights. The Reactive
Controller uses the top retrieved summaries as hints rather than executable
commands.

When the Event Trigger escalates control, the Strategic Controller
consults SA-KG for plan checking rather than local traces alone.
SA-KG stores validated transitions as state nodes, action edges, and similarity
edges, inspired by state-action-result graphs in experience-driven agents
without privileged environment APIs~\citep{tang2025experiencedriven} but using
retrieved graph fragments as evidence rather than directly executing retrieved
actions.
Starting from the nearest state nodes for the current subgoal, SA-KG retrieves
graph fragments by vector search, keeps matches above $\tau_{\mathrm{KG}}$, and
ranks fragment $j$ by
\begin{equation}
\label{eq:sakg_score}
S_{\mathrm{KG}}(j \mid x_t)=\beta_C C_j+\beta_P P_j,
\end{equation}
where $C_j=1-\mathrm{dist}(x_t,s_j)$ is vector similarity and $P_j$ is the
success rate of the best outgoing action, with execution count used to break
ties. Fusion first min--max normalizes $S_{\mathrm{MB}}$ and $S_{\mathrm{KG}}$
within the current candidate sets, then uses
$S_{\mathrm{fuse}}=\lambda_{\mathrm{MB}}\tilde S_{\mathrm{MB}}+\lambda_{\mathrm{KG}}\tilde S_{\mathrm{KG}}$;
unmatched graph actions receive zero normalized SA-KG score. After execution,
\method{} writes only successful execution, interpretable recovery, or
planner-validated progress, preventing transient errors from becoming long-term
evidence. Appendix Sec.~\mbox{\ref{sec:appendix_memory_details}} gives the
storage fields, fusion score, fixed weights, and promotion rules.
